% ---------------------------------------------------------------------------
% INSERT FILE: cutoff_cost_cite_insert.tex   (NS-3, drafted 2026-09-21)
% purpose: place prop:v125-cutoff-cost in the published doubly-exponential
%          family (Zhu's frequency threshold, Groskin's brute-cutoff cost) and
%          state that the objects differ.  Source: COMPARISON.md.
% anchor:  manuscript/fixed_space_prime_action_v1.tex, immediately after the
%          statement of prop:v125-cutoff-cost.  Its four-line proof follows the
%          statement directly, so the insert goes after that proof's
%            \end{proof}
%          i.e. immediately BEFORE the line
%            For $c=10^{-8}$, freshly evaluating the physical weighted prime bound
%          (placing prose between a proposition and its proof is avoided).
% cites:   \cite{Zhu2026} (present), \cite{Groskin2026} (see bib_inserts.tex).
% status:  draft for integration; not yet in the manuscript
% ---------------------------------------------------------------------------

The doubly exponential shape of \eqref{eq:v125-cutoff-cost}, namely
$N\sim\exp(e^{a})$ in the half-width $a=\log\lambda$, is that of a
published family: Zhu~\cite{Zhu2026} shows that any one-stroke
certificate must resolve frequencies up to $2\pi e^{A_L}$ with
$A_L\sim4e^{L}$ in the half-width $L$, drawing the same conclusion that
the positivity route cannot reach RH unassisted, and
Groskin~\cite{Groskin2026} gives the corresponding brute-cutoff cost
for the Connes--van Suijlekom truncations (an archimedean cutoff of
order $10^{63}$ to resolve a spectral scale $10^{-59}$ at prime cutoff
$c=100$), resolved there instead by a cutoff-free interval $LDL^{T}$
factorization, the factorization technique used here.
The objects differ: Zhu bounds the frequency resolution of a
one-stroke certificate, Groskin the archimedean cutoff of a Galerkin
truncation, and the proposition above the remote cutoff of one specific
scalar majorant; the proposition is placed in that family and is not
presented as new.
